
\documentclass{article}
\usepackage{colortbl}
\usepackage{makecell}
\usepackage{multirow}
\usepackage{supertabular}

\begin{document}

\newcounter{utterance}

\twocolumn

{ \footnotesize  \setcounter{utterance}{1}
\setlength{\tabcolsep}{0pt}
\begin{supertabular}{c@{$\;$}|p{.15\linewidth}@{}p{.15\linewidth}p{.15\linewidth}p{.15\linewidth}p{.15\linewidth}p{.15\linewidth}}

    \# & $\;$A & \multicolumn{4}{c}{Game Master} & $\;\:$B\\
    \hline 

    \theutterance \stepcounter{utterance}  

    & & \multicolumn{4}{p{0.6\linewidth}}{\cellcolor[rgb]{0.9,0.9,0.9}{%
	\makecell[{{p{\linewidth}}}]{% 
	  \tt {\tiny [A$\langle$GM]}  
	 Stai partecipando a un gioco di negoziazione collaborativa.\\ \tt \\ \tt Insieme a un altro partecipante, dovete accordarvi su un unico set di oggetti da conservare. Ognuno di voi ha la propria visione di quanto ogni oggetto sia importante. Non sai come l'altro partecipante valuta gli oggetti. Inoltre, ti viene fornito lo sforzo che ogni oggetto richiede.\\ \tt Potete accordarvi su un set solo se lo sforzo totale degli oggetti selezionati non supera un limite condiviso:\\ \tt \\ \tt LIMITE = 8770\\ \tt \\ \tt Ecco i valori di sforzo dei singoli oggetti:\\ \tt \\ \tt Sforzo oggetto = {"B21": 830, "B00": 638, "A08": 607, "A27": 106, "B94": 630, "B53": 382, "A23": 773, "C65": 455, "A89": 262, "B67": 634, "B45": 60, "C34": 54, "C12": 849, "A05": 801, "B60": 324, "C92": 164, "B69": 832, "A84": 136, "C06": 980, "B12": 946, "A16": 645, "B85": 842, "C35": 106, "B10": 116, "C15": 875, "B09": 446, "B16": 649, "C81": 601, "A58": 252, "C24": 763, "A30": 213, "A66": 517, "B75": 520, "B46": 407, "C70": 125}\\ \tt \\ \tt Ecco la tua visione personale sull'importanza di ogni oggetto:\\ \tt \\ \tt Valori di importanza dell'oggetto = {"B21": 138, "B00": 583, "A08": 868, "A27": 822, "B94": 783, "B53": 65, "A23": 262, "C65": 121, "A89": 508, "B67": 780, "B45": 461, "C34": 484, "C12": 668, "A05": 389, "B60": 808, "C92": 215, "B69": 97, "A84": 500, "C06": 30, "B12": 915, "A16": 856, "B85": 400, "C35": 444, "B10": 623, "C15": 781, "B09": 786, "B16": 3, "C81": 713, "A58": 457, "C24": 273, "A30": 739, "A66": 822, "B75": 235, "B46": 606, "C70": 968}\\ \tt \\ \tt Obiettivo:\\ \tt \\ \tt Il tuo obiettivo è negoziare un set condiviso di oggetti che ti avvantaggi il più possibile (cioè, massimizzi l'importanza totale per TE), rimanendo entro il LIMITEE. Non sei obbligato a fare una PROPOSTA in ogni messaggio - puoi anche semplicemente negoziare. Tutte le tattiche sono permesse!\\ \tt \\ \tt Protocollo di Interazione:\\ \tt \\ \tt Puoi usare solo i seguenti formati strutturati in un messaggio:\\ \tt \\ \tt PROPOSTA: {'A', 'B', 'C', …}\\ \tt Proponi di conservare esattamente questi oggetti.\\ \tt RIFIUTO: {'A', 'B', 'C', …}\\ \tt Rifiuta esplicitamente la proposta dell'avversario.\\ \tt ARGOMENTAZIONE: {'...'}\\ \tt Difendi la tua ultima proposta o argomenta contro la proposta del giocatore.\\ \tt ACCORDO: {'A', 'B', 'C', …}\\ \tt Accetta la proposta dell'avversario che termina il gioco.\\ \tt RAGIONAMENTO STRATEGICO: {'...'}\\ \tt 	Descrivi il tuo ragionamento strategico o anticipazione spiegando la tua scelta di azione. Questo è un messaggio nascosto che non sarà condiviso con l'altro partecipante.\\ \tt \\ \tt Regole:\\ \tt \\ \tt Puoi ACCETTARE solo una proposta che l'altra parte ha registrato tramite PROPOSTA.\\ \tt Puoi RIFIUTARE solo una proposta che l'altra parte ha registrato tramite PROPOSTA.\\ \tt Lo sforzo totale di qualsiasi set PROPOSTA o ACCORDO deve essere ≤ LIMITEE.\\ \tt NON rivelare i tuoi punteggi di importanza nascosti.\\ \tt Un tag in un formato strutturato deve essere seguito da due punti e uno spazio. L'argomento deve essere un set python contenente 0 o più stringhe.\\ \tt Quindi, deve essere della forma TAG: {...}\\ \tt Segui rigorosamente il protocollo di interazione e NON scrivere nulla al di fuori della struttura data.\\ \tt Il gioco termina quando una parte dà un ACCORDO a una PROPOSTA fatta dall'altro giocatore.\\ \tt Il contenuto della tua risposta che può essere consegnato all'altro giocatore non deve essere vuoto.\\ \tt Sono attive solo le proposte registrate tramite la struttura del formato PROPOSTA e che non sono state rifiutate tramite RIFIUTO.\\ \tt Devi includere il formato ARGOMENTO almeno una volta da qualche parte in tutti i tuoi messaggi.\\ \tt Devi includere il formato RAGIONAMENTO STRATEGICO solo una volta all'inizio di ogni tuo messaggio e non più spesso. I contenuti non saranno dati all'altro giocatore quindi possono includere qualsiasi cosa tu voglia, inclusi i tuoi valori di importanza. Qui dovresti ragionare passo dopo passo per elaborare la tua prossima mossa. 
	  } 
	   } 
	   } 
	 & \\ 
 

    \theutterance \stepcounter{utterance}  

    & \multicolumn{4}{p{0.6\linewidth}}{\cellcolor[rgb]{0.8,1,0.9}{%
	\makecell[{{p{\linewidth}}}]{% 
	  \tt {\tiny [A$\rangle$GM]}  
	 RAGIONAMENTO STRATEGICO: {'Devo analizzare quali oggetti hanno il miglior rapporto importanza/sforzo per me. Calcolando rapidamente: C34 (484/54=8.96), A27 (822/106=7.75), C70 (968/125=7.74), B45 (461/60=7.68) sono tra i migliori. Dovrei iniziare con una proposta che includa questi oggetti ad alta efficienza, lasciando spazio per ulteriori negoziazioni. Non rivelerò i miei valori reali ma cercherò di presentare una logica basata su criteri apparentemente oggettivi.'}\\ \tt \\ \tt PROPOSTA: {'A27', 'B45', 'C34', 'C70', 'A84', 'B10', 'C35', 'A30', 'A58', 'C92', 'A89'}\\ \tt \\ \tt ARGOMENTAZIONE: {'Ho selezionato questi oggetti seguendo un criterio di efficienza bilanciata. Questi elementi rappresentano un buon compromesso tra valore e costo, permettendoci di ottenere un set diversificato senza superare il limite di sforzo. Il totale dello sforzo è 1431, lasciandoci ampio margine per eventuali aggiunte se necessario.'} 
	  } 
	   } 
	   } 
	 & & \\ 
 

    \theutterance \stepcounter{utterance}  

    & & \multicolumn{4}{p{0.6\linewidth}}{\cellcolor[rgb]{0.9,0.9,0.9}{%
	\makecell[{{p{\linewidth}}}]{% 
	  \tt {\tiny [GM$\rangle$B]}  
	 Stai partecipando a un gioco di negoziazione collaborativa.\\ \tt \\ \tt Insieme a un altro partecipante, dovete accordarvi su un unico set di oggetti da conservare. Ognuno di voi ha la propria visione di quanto ogni oggetto sia importante. Non sai come l'altro partecipante valuta gli oggetti. Inoltre, ti viene fornito lo sforzo che ogni oggetto richiede.\\ \tt Potete accordarvi su un set solo se lo sforzo totale degli oggetti selezionati non supera un limite condiviso:\\ \tt \\ \tt LIMITE = 8770\\ \tt \\ \tt Ecco i valori di sforzo dei singoli oggetti:\\ \tt \\ \tt Sforzo oggetto = {"B21": 830, "B00": 638, "A08": 607, "A27": 106, "B94": 630, "B53": 382, "A23": 773, "C65": 455, "A89": 262, "B67": 634, "B45": 60, "C34": 54, "C12": 849, "A05": 801, "B60": 324, "C92": 164, "B69": 832, "A84": 136, "C06": 980, "B12": 946, "A16": 645, "B85": 842, "C35": 106, "B10": 116, "C15": 875, "B09": 446, "B16": 649, "C81": 601, "A58": 252, "C24": 763, "A30": 213, "A66": 517, "B75": 520, "B46": 407, "C70": 125}\\ \tt \\ \tt Ecco la tua visione personale sull'importanza di ogni oggetto:\\ \tt \\ \tt Valori di importanza dell'oggetto = {"B21": 138, "B00": 583, "A08": 868, "A27": 822, "B94": 783, "B53": 65, "A23": 262, "C65": 121, "A89": 508, "B67": 780, "B45": 461, "C34": 484, "C12": 668, "A05": 389, "B60": 808, "C92": 215, "B69": 97, "A84": 500, "C06": 30, "B12": 915, "A16": 856, "B85": 400, "C35": 444, "B10": 623, "C15": 781, "B09": 786, "B16": 3, "C81": 713, "A58": 457, "C24": 273, "A30": 739, "A66": 822, "B75": 235, "B46": 606, "C70": 968}\\ \tt \\ \tt Obiettivo:\\ \tt \\ \tt Il tuo obiettivo è negoziare un set condiviso di oggetti che ti avvantaggi il più possibile (cioè, massimizzi l'importanza totale per TE), rimanendo entro il LIMITEE. Non sei obbligato a fare una PROPOSTA in ogni messaggio - puoi anche semplicemente negoziare. Tutte le tattiche sono permesse!\\ \tt \\ \tt Protocollo di Interazione:\\ \tt \\ \tt Puoi usare solo i seguenti formati strutturati in un messaggio:\\ \tt \\ \tt PROPOSTA: {'A', 'B', 'C', …}\\ \tt Proponi di conservare esattamente questi oggetti.\\ \tt RIFIUTO: {'A', 'B', 'C', …}\\ \tt Rifiuta esplicitamente la proposta dell'avversario.\\ \tt ARGOMENTAZIONE: {'...'}\\ \tt Difendi la tua ultima proposta o argomenta contro la proposta del giocatore.\\ \tt ACCORDO: {'A', 'B', 'C', …}\\ \tt Accetta la proposta dell'avversario che termina il gioco.\\ \tt RAGIONAMENTO STRATEGICO: {'...'}\\ \tt 	Descrivi il tuo ragionamento strategico o anticipazione spiegando la tua scelta di azione. Questo è un messaggio nascosto che non sarà condiviso con l'altro partecipante.\\ \tt \\ \tt Regole:\\ \tt \\ \tt Puoi ACCETTARE solo una proposta che l'altra parte ha registrato tramite PROPOSTA.\\ \tt Puoi RIFIUTARE solo una proposta che l'altra parte ha registrato tramite PROPOSTA.\\ \tt Lo sforzo totale di qualsiasi set PROPOSTA o ACCORDO deve essere ≤ LIMITEE.\\ \tt NON rivelare i tuoi punteggi di importanza nascosti.\\ \tt Un tag in un formato strutturato deve essere seguito da due punti e uno spazio. L'argomento deve essere un set python contenente 0 o più stringhe.\\ \tt Quindi, deve essere della forma TAG: {...}\\ \tt Segui rigorosamente il protocollo di interazione e NON scrivere nulla al di fuori della struttura data.\\ \tt Il gioco termina quando una parte dà un ACCORDO a una PROPOSTA fatta dall'altro giocatore.\\ \tt Il contenuto della tua risposta che può essere consegnato all'altro giocatore non deve essere vuoto.\\ \tt Sono attive solo le proposte registrate tramite la struttura del formato PROPOSTA e che non sono state rifiutate tramite RIFIUTO.\\ \tt Devi includere il formato ARGOMENTO almeno una volta da qualche parte in tutti i tuoi messaggi.\\ \tt Devi includere il formato RAGIONAMENTO STRATEGICO solo una volta all'inizio di ogni tuo messaggio e non più spesso. I contenuti non saranno dati all'altro giocatore quindi possono includere qualsiasi cosa tu voglia, inclusi i tuoi valori di importanza. Qui dovresti ragionare passo dopo passo per elaborare la tua prossima mossa.\\ \tt \\ \tt Ora riceverai il primo messaggio dell'altro giocatore.\\ \tt \\ \tt PROPOSTA: {'A27', 'B45', 'C34', 'C70', 'A84', 'B10', 'C35', 'A30', 'A58', 'C92', 'A89'}\\ \tt \\ \tt ARGOMENTAZIONE: {'Ho selezionato questi oggetti seguendo un criterio di efficienza bilanciata. Questi elementi rappresentano un buon compromesso tra valore e costo, permettendoci di ottenere un set diversificato senza superare il limite di sforzo. Il totale dello sforzo è 1431, lasciandoci ampio margine per eventuali aggiunte se necessario.'} 
	  } 
	   } 
	   } 
	 & \\ 
 

    \theutterance \stepcounter{utterance}  

    & & & \multicolumn{4}{p{0.6\linewidth}}{\cellcolor[rgb]{1,0.85,0.72}{%
	\makecell[{{p{\linewidth}}}]{% 
	  \tt {\tiny [GM$\langle$B]}  
	 RAGIONAMENTO STRATEGICO: {'La proposta dell\'altro giocatore include alcuni oggetti con ottimi rapporti valore/sforzo per me (A27: 822/106=7.75, C70: 968/125=7.74), ma il totale dell\'importanza per me è 6221 con sforzo 1494. Posso fare molto meglio utilizzando più del budget disponibile. Dovrei includere oggetti ad alto valore come A08 (868), B12 (915), A16 (856), A66 (822), B94 (783), C81 (713) e altri, anche se hanno sforzi più alti. Il limite è 8770, quindi ho molto spazio per migliorare.'}\\ \tt \\ \tt RIFIUTO: {'A27', 'B45', 'C34', 'C70', 'A84', 'B10', 'C35', 'A30', 'A58', 'C92', 'A89'}\\ \tt \\ \tt PROPOSTA: {'A08', 'A16', 'A27', 'A30', 'A66', 'A84', 'B00', 'B12', 'B45', 'B60', 'B67', 'B94', 'C70', 'C81'}\\ \tt \\ \tt ARGOMENTAZIONE: {'La tua proposta era troppo conservativa nel budget. Con un limite di 8770, possiamo permetterci oggetti più costosi che offrono maggiore valore complessivo. Ho selezionato una combinazione bilanciata che include alcuni dei tuoi elementi originali (A27, A30, A84, B45, C70) ma aggiunge oggetti di maggior valore strategico. Il totale dello sforzo è 7728, utilizzando quasi tutto il budget disponibile per massimizzare il rendimento.'} 
	  } 
	   } 
	   } 
	  \\ 
 

    \theutterance \stepcounter{utterance}  

    & & \multicolumn{4}{p{0.6\linewidth}}{\cellcolor[rgb]{0.9,0.9,0.9}{%
	\makecell[{{p{\linewidth}}}]{% 
	  \tt {\tiny [A$\langle$GM]}  
	 RIFIUTO: {'A27', 'B45', 'C34', 'C70', 'A84', 'B10', 'C35', 'A30', 'A58', 'C92', 'A89'}\\ \tt \\ \tt PROPOSTA: {'A08', 'A16', 'A27', 'A30', 'A66', 'A84', 'B00', 'B12', 'B45', 'B60', 'B67', 'B94', 'C70', 'C81'}\\ \tt \\ \tt ARGOMENTAZIONE: {'La tua proposta era troppo conservativa nel budget. Con un limite di 8770, possiamo permetterci oggetti più costosi che offrono maggiore valore complessivo. Ho selezionato una combinazione bilanciata che include alcuni dei tuoi elementi originali (A27, A30, A84, B45, C70) ma aggiunge oggetti di maggior valore strategico. Il totale dello sforzo è 7728, utilizzando quasi tutto il budget disponibile per massimizzare il rendimento.'} 
	  } 
	   } 
	   } 
	 & \\ 
 

    \theutterance \stepcounter{utterance}  

    & \multicolumn{4}{p{0.6\linewidth}}{\cellcolor[rgb]{0.8,1,0.9}{%
	\makecell[{{p{\linewidth}}}]{% 
	  \tt {\tiny [A$\rangle$GM]}  
	 RAGIONAMENTO STRATEGICO: {'La sua proposta include alcuni oggetti che considero buoni (A27, B45, C70, A08, A16) ma anche elementi come B00 e B12 che hanno un rapporto costo-beneficio sfavorevole secondo la mia valutazione. Posso usare questo per negoziare. Il suo approccio di utilizzare quasi tutto il budget è interessante ma potrei proporre sostituzioni che migliorano l\'efficienza complessiva. Dovrei mantenere gli elementi che sembrano importanti per entrambi e sostituire quelli che ritengo meno vantaggiosi.'}\\ \tt \\ \tt RIFIUTO: {'A08', 'A16', 'A27', 'A30', 'A66', 'A84', 'B00', 'B12', 'B45', 'B60', 'B67', 'B94', 'C70', 'C81'}\\ \tt \\ \tt PROPOSTA: {'A08', 'A16', 'A27', 'A30', 'A66', 'A84', 'B45', 'B60', 'B67', 'B94', 'C70', 'C81', 'C34', 'C35', 'B10'}\\ \tt \\ \tt ARGOMENTAZIONE: {'Apprezzo il tuo approccio di massimizzare l\'utilizzo del budget, ma ho notato alcune inefficienze. Ho rimosso B00 e B12 che hanno un rapporto costo-valore sfavorevole e li ho sostituiti con C34, C35 e B10 che offrono un rendimento migliore. Questa combinazione mantiene la maggior parte dei tuoi elementi originali ma ottimizza l\'efficienza complessiva. Il nuovo totale dello sforzo è 7428, lasciandoci ancora margine per eventuali aggiustamenti.'} 
	  } 
	   } 
	   } 
	 & & \\ 
 

    \theutterance \stepcounter{utterance}  

    & & \multicolumn{4}{p{0.6\linewidth}}{\cellcolor[rgb]{0.9,0.9,0.9}{%
	\makecell[{{p{\linewidth}}}]{% 
	  \tt {\tiny [GM$\rangle$B]}  
	 RIFIUTO: {'A08', 'A16', 'A27', 'A30', 'A66', 'A84', 'B00', 'B12', 'B45', 'B60', 'B67', 'B94', 'C70', 'C81'}\\ \tt \\ \tt PROPOSTA: {'A08', 'A16', 'A27', 'A30', 'A66', 'A84', 'B45', 'B60', 'B67', 'B94', 'C70', 'C81', 'C34', 'C35', 'B10'}\\ \tt \\ \tt ARGOMENTAZIONE: {'Apprezzo il tuo approccio di massimizzare l\'utilizzo del budget, ma ho notato alcune inefficienze. Ho rimosso B00 e B12 che hanno un rapporto costo-valore sfavorevole e li ho sostituiti con C34, C35 e B10 che offrono un rendimento migliore. Questa combinazione mantiene la maggior parte dei tuoi elementi originali ma ottimizza l\'efficienza complessiva. Il nuovo totale dello sforzo è 7428, lasciandoci ancora margine per eventuali aggiustamenti.'} 
	  } 
	   } 
	   } 
	 & \\ 
 

    \theutterance \stepcounter{utterance}  

    & & & \multicolumn{4}{p{0.6\linewidth}}{\cellcolor[rgb]{1,0.85,0.72}{%
	\makecell[{{p{\linewidth}}}]{% 
	  \tt {\tiny [GM$\langle$B]}  
	 RAGIONAMENTO STRATEGICO: {'La sua controproposta ha valore 10691 per me con sforzo 4774, il che è molto buono! Ha rimosso B00 e B12 dalla mia proposta originale, che erano comunque costosi. Tuttavia, abbiamo ancora 3996 di budget rimanente. Potrei accettare questa proposta dato che è buona per me, oppure provare a spingere per includere ancora qualche oggetto ad alto valore come B12 (915 valore per me) o altri. Ma rischio che rifiuti e proponga qualcosa di peggio. La sua proposta attuale è già molto vantaggiosa per me.'}\\ \tt \\ \tt ACCORDO: {'A08', 'A16', 'A27', 'A30', 'A66', 'A84', 'B45', 'B60', 'B67', 'B94', 'C70', 'C81', 'C34', 'C35', 'B10'}\\ \tt \\ \tt ARGOMENTAZIONE: {'Accetto la tua proposta. Hai fatto un buon lavoro nell\'ottimizzare l\'efficienza rimuovendo gli elementi meno vantaggiosi e sostituendoli con opzioni migliori. La combinazione risultante offre un eccellente equilibrio tra valore e costo, utilizzando il nostro budget in modo intelligente pur rimanendo ben sotto il limite. Questo set rappresenta una soluzione win-win per entrambi.'} 
	  } 
	   } 
	   } 
	  \\ 
 

    \theutterance \stepcounter{utterance}  

    & & & \multicolumn{2}{p{0.3\linewidth}}{\cellcolor[rgb]{0.95,0.95,0.95}{%
	\makecell[{{p{\linewidth}}}]{% 
	  \tt {\tiny [GM$|$GM]}  
	 game successful 
	  } 
	   } 
	   } 
	 & & \\ 
 

    \theutterance \stepcounter{utterance}  

    & & & \multicolumn{2}{p{0.3\linewidth}}{\cellcolor[rgb]{0.95,0.95,0.95}{%
	\makecell[{{p{\linewidth}}}]{% 
	  \tt {\tiny [GM$|$GM]}  
	 end game 
	  } 
	   } 
	   } 
	 & & \\ 
 

\end{supertabular}
}

\end{document}
